\section{The sign, the evanescent part, and the Hodge-index route}

The sign-only construction of SG-08--SG-11 proves the required finite-field absolute values. Its decisive sign is the parity of an \emph{effective Frobenius remainder}: generic weak Lefschetz makes a Gysin map surjective, and the remaining virtual class is the class of an actual kernel together with a tensor power. This does not identify that sign with positive definiteness of an intersection form. Both arguments use hyperplane sections and evanescent cohomology, but the effective-remainder argument avoids the positive involution on correspondences required by the index route.

The supplied cote~16 scans~034, 036, 038 and~040 describe restriction, Gysin, Lefschetz operators, the evanescent summand, and its projector. They do not state positivity on the whole evanescent space. Scan~049 explicitly distinguishes algebraicity of an inverse from the extra index relation $I(X\times X)$ used for splittings and semisimplicity. Scan~050 announces a deduction of Weil from $C$ and $A$. The earlier definitions of those two labels are absent from the supplied excerpts, so their identification with modern standardized letters is not silently assumed. The conclusion here is precise: the sign-only proof realizes a cohomological hyperplane-induction programme, but it does not prove the additional cycle-theoretic hypotheses of the positive-involution argument announced on scans~049--050. The unrelated printed material on scans~035, 037 and~039 is an inventory fact, not by itself a transcription error.

\begin{proposition}[SG-01: parity is not intersection signature]
The sign $(-1)^j$ of a cohomology group in a trace formula, or its exponent $(-1)^{j+1}$ in the cohomological zeta function, does not determine the raw intersection sign on primitive algebraic classes in that degree. Neither sign gives a definite intersection form on the whole evanescent part of a hyperplane section.
\end{proposition}

\begin{proof}
Let $Q^{2r}$ be a split smooth even-dimensional quadric, as a smooth hyperplane section of a split $Q^{2r+1}$. Let $a,b$ be its two classes of maximal linear $r$-planes and put $\zeta=a-b$. The explicit calculation in SG-02 gives
\[
 h\zeta=0,\qquad \deg(\zeta^2)=2(-1)^r,
 \qquad E^{2r}\simeq\mathbb Q_\ell(-r).
\]
Here $E^{2r}$ is the kernel of the Gysin map
\[
 H^{2r}(Q^{2r}_{\overline{\mathbb F}_q},\mathbb Q_\ell)
 \longrightarrow H^{2r+2}(Q^{2r+1}_{\overline{\mathbb F}_q},\mathbb Q_\ell(1)).
\]
Its trace contribution over $\mathbb F_{q^n}$ is $+q^{rn}$, and its determinant occurs with exponent $-1$ in the zeta function. These two signs are independent of $r$, whereas $\zeta^2$ alternates with $r$. Thus $Q^2,Q^4,Q^6$ have trace signs $+,+,+$, determinant exponents $-,-,-$, and primitive squares $-2,+2,-2$. A uniform identification of parity with the raw index sign is false.

After normalizing the primitive form to $(-1)^r\deg(xy)$, all three Tate lines are positive. That numerical agreement with a positive multiplicity does not identify the constructions: one is a bilinear inequality on rational algebraic classes; the other is a sign in the Grothendieck group of Frobenius modules, defined without a pairing.

There is also no definite intersection form on the entire evanescent space in general. Embed $X=\mathbb P^2_{\mathbb C}$ by $\mathcal O(3)$, and take a smooth plane cubic $Y$, now a hyperplane section. Since $H^1(X)=H^3(X)=0$, its evanescent part is $E^1=H^1(Y,\mathbb Q)$, of dimension two. Its cup-product pairing is alternating by graded commutativity, and is nondegenerate by duality. Every vector has square zero; this cannot be a positive- or negative-definite real symmetric form. Its cohomological trace sign is nevertheless uniformly negative. This example concerns the whole evanescent space and does not contradict the Hodge standard conjecture, which concerns primitive \emph{algebraic} classes in even degree with a prescribed normalization.
\end{proof}

The exact cohomological locators are \cite[eq.~(1.5.4), p.~276]{WeilI} for determinant parity, \cite[\S4.1, p.~287]{WeilI} for the local vanishing-cycle squares $2,0,-2,0$ in dimensions congruent to $0,1,2,3$ modulo~$4$, and \cite[\S5.9, p.~293]{WeilI} for the alternating or symmetric form on $E/(E\cap E^\perp)$. The local Picard--Lefschetz signs are genuine intersection signs. They still do not assert positive definiteness on all of $E$.

\subsection{The split sixfold quadric and its Tate factor}

\begin{proposition}[SG-02: the split-quadric calculation]
Let $r\ge1$. On a split smooth $Q^{2r}$, the primitive middle cohomology and its corresponding rational Chow-motive summand are respectively $\mathbb Q_\ell(-r)$ and $\mathbb Q(-r)$. If $\zeta$ is the difference of the two maximal-plane classes, then $\deg(\zeta^2)=2(-1)^r$. Moreover,
\begin{align*}
 \#Q^{2r}(\mathbb F_{q^n})&=\sum_{j=0}^{2r}q^{jn}+q^{rn},\\
 Z(Q^{2r}/\mathbb F_q,T)&=
 \left((1-q^rT)\prod_{j=0}^{2r}(1-q^jT)\right)^{-1}.
\end{align*}
For the split sixfold model over $\mathbb Z$, the primitive global factor is $\zeta(s-3)$. Under the stipulated Alp\"oge Riemann hypothesis its nontrivial zeros have real part $7/2$. This is distinct from the local statement that its degree-six Frobenius eigenvalue is $q^3$.
\end{proposition}

\begin{proof}
Write the equation as $\sum_{i=0}^r x_i y_i=0$. The chart $x_0\ne0$ is $\mathbb A^{2r}$. Its complement is the projective cone on $Q^{2r-2}$: it has one vertex, and its complement of the vertex is an $\mathbb A^1$-bundle over $Q^{2r-2}$. Starting with $Q^0$, which consists of two points, this gives an affine paving with one cell in every dimension $0,1,\ldots,2r$ and an extra cell in dimension $r$. All these cells are defined over the split base. Counting the cells proves the point-count formula, and exponentiating $\sum_{n\ge1}\#Q(\mathbb F_{q^n})T^n/n$ proves the zeta formula. The cell filtration and the localization sequence give only even cohomology, with one Tate class in every even degree and an additional middle class. The two maximal-plane classes below form an explicit Frobenius eigenbasis in middle degree, so no semisimplicity or weight theorem is being used to split a possible extension.

For a maximal plane $\mathbb P^r\subset Q^{2r}$ the normal exact sequence is
\[
 0\longrightarrow N_{\mathbb P^r/Q^{2r}}
 \longrightarrow\mathcal O(1)^{r+1}
 \longrightarrow\mathcal O(2)\longrightarrow0.
\]
Consequently its self-intersection is the coefficient of $t^r$ in
$(1+t)^{r+1}/(1+2t)$, namely $(1+(-1)^r)/2$. A special linear
$\mathbb P^{r+1}$ section is the union of a plane from each ruling, so
$h^r=a+b$. Since $\deg(h^r a)=\deg(h^r b)=1$, it follows that
\[
 a^2=b^2=\frac{1+(-1)^r}{2},\qquad
 ab=\frac{1-(-1)^r}{2},\qquad
 (a-b)^2=2(-1)^r.
\]
The intersection matrix is nonsingular, proving the independence of the two classes. In codimension $r+1$ the cell calculation gives a single generator, and $ha=hb$ by pairing with $h^{r-1}$; hence $h\zeta=0$. The same equality holds in the Chow ring: the cell closures generate its groups, and their cycle classes, of the same rank, are independent. The correspondence
\[
 p_{\mathrm{prim}}=\frac{\zeta\times\zeta}{\deg(\zeta^2)}
\]
acts identically on $\zeta$ and kills the hyperplane line. The cycle $\zeta$ and its transposed cycle divided by $\deg(\zeta^2)$ give inverse maps between this summand and $\mathbb Q(-r)$. This proves the assertion about the rational Chow motive. For $Q^{2r}\subset Q^{2r+1}$, the ambient middle class has nonzero Gysin image and $\zeta$ has zero Gysin image. Since the kernel has dimension one, it is exactly this primitive Tate line.

For $r=3$, the model $\sum_{i=0}^3x_i y_i=0$ in $\mathbb P^7_{\mathbb Z}$ is smooth over $\mathbb Z$: its partial derivatives are the coordinates $x_i,y_i$, which cannot vanish simultaneously at a projective point. Every finite-field fibre is split. Its primitive Euler factor is $(1-p^{3-s})^{-1}$, so their product is $\zeta(s-3)$ in the half-plane of absolute convergence, and then by meromorphic continuation. The change of variable $u=s-3$ carries the stipulated Riemann hypothesis to $\operatorname{Re}(s)=7/2$.
\end{proof}

Thus the requested absolute values for the ambient quadric have a direct Tate-motive computation; SG-08--SG-11 also give them through the general sign-only theorem. The elementary intersection inequality is independently visible as $(-1)^r\zeta^2=2>0$. Neither computation should be confused with a deduction of the Hodge standard conjecture from trace parity.

\subsection{The extra datum in the positive-involution proof}

\begin{lemma}[SG-03: the exact positivity mechanism]
Let $A$ be a finite-dimensional real algebra with identity, $\dagger$ an involutive anti-automorphism, and $\tau:A\to\mathbb R$ a cyclic linear functional satisfying $\tau(a^\dagger)=\tau(a)$. Suppose
\[
 B(a,b)=\tau(ab^\dagger)
\]
is positive definite. If $F^\dagger F=c\,1$ with $c>0$, every eigenvalue of $F$ in a faithful finite-dimensional representation has modulus $\sqrt c$.
\end{lemma}

\begin{proof}
The equation makes $F$ invertible, hence $FF^\dagger=c$ also. Cyclicity gives
\[
 B(Fa,Fb)=\tau(Fab^\dagger F^\dagger)
        =\tau(F^\dagger F ab^\dagger)=cB(a,b).
\]
Left multiplication by $c^{-1/2}F$ is therefore an isometry of a positive-definite real vector space. Its complex eigenvalues have modulus one. The minimum polynomial of left multiplication by $F$ equals the minimum polynomial of $F$ as an element of $A$: indeed $p(L_F)=L_{p(F)}$, and applying to $1$ detects zero. A faithful representation has that same minimum polynomial, proving the assertion.
\end{proof}

This short lemma states the precise form needed here of the classical positive-Rosati argument; compare \cite[Th.~13.2 and Lemmas~14.2--14.3]{MilneAV}. In the correspondence-algebra route, construction of the positive form and its algebraic adjoint is additional content. Merely knowing an effective trace remainder does not produce it. For a sharp formal comparison, the indefinite matrix
\[
 J=\begin{pmatrix}0&1\\1&0\end{pmatrix},\qquad
 F=\operatorname{diag}(\sqrt c\,a,\sqrt c/a),\quad a>1,
\]
satisfies $F^tJF=cJ$ but has eigenvalues of different absolute values. Conversely $F=\sqrt c\,I$ has pure absolute values while preserving this indefinite form up to the same multiplier. These are countermodels between linear algebra data, not geometric counterexamples to the Hodge standard conjecture. Cote~16 scan~049 is precisely where the extra index input is acknowledged.

\subsection{CM elliptic powers and arbitrary abelian varieties}

\begin{proposition}[SG-04: the precise hexagonal-curve case]
Let $E/\mathbb C$ be the elliptic curve $\mathbb C/\mathbb Z[\omega]$. Every rational algebraic cohomology class on $E^n$ is generated by divisor classes. The same assertion holds for the geometric rational algebraic $\ell$-adic classes on powers of every good finite-field reduction of $E$. The Hodge standard conjecture on these class spaces is consequently an unconditional instance of \cite[Remarks~3.7--3.8]{MilnePolarizations}. This concerns cohomology classes, and makes no assertion that every Chow class or every class modulo algebraic equivalence is a polynomial in divisors.
\end{proposition}

\begin{proof}
Over $\mathbb C$, put $K=\mathbb Q(\omega)$. Then
$\operatorname{End}^0(E^n)=M_n(K)$, and Rosati for the product polarization is conjugate transpose. Its fixed subspace has rational dimension
$n+2\binom n2=n^2$. The N\'eron--Severi/Rosati identification
\cite[Prop.~13.1]{MilneAV} therefore gives $\dim_{\mathbb Q}\operatorname{NS}(E^n)_{\mathbb Q}=n^2$. The first Chern class is injective on this space. Since $H^{1,1}(E^n)$ has complex dimension $n^2$, divisor classes span it over $\mathbb C$.
Write $H^1(E^n,\mathbb C)=U\oplus W$ in Hodge types $(1,0),(0,1)$, each of dimension $n$. The space
\[
 H^{r,r}(E^n)=\bigwedge^rU\otimes\bigwedge^rW
\]
is spanned by products $(u_1\wedge w_1)\cdots(u_r\wedge w_r)$ of degree-two classes. Thus the complexification of the rational algebra generated by divisors fills $H^{r,r}$. Every rational algebraic class belongs to that rational algebra. Here we use the exterior-algebra description of abelian cohomology, \cite[Th.~11.1(b), and Th.~2.3 for Betti cohomology]{MilneAV}.

For completeness the finite-field assertion requires both geometric endomorphisms and rational descent of coefficients. The ordinary/supersingular classification and the criterion $p\nmid\operatorname{Tr}(\pi_E)$ in the ordinary case are \cite[Th.~13.3, Th.~13.18, Cor.~13.20]{SutherlandEC}. For an ordinary elliptic curve, the two roots $\alpha,\beta$ of the Frobenius polynomial satisfy $\alpha\beta=q$ and $\alpha/\beta$ is not a root of unity. Indeed they are algebraic integers, their sum is a $p$-adic unit and their product has positive valuation, so one has valuation zero and the other has valuation $v_p(q)$. A geometric algebraic class is fixed after a finite extension and a Tate twist. In degree $2r$ on $E^n$, the invariant monomials therefore have exactly $r$ vectors from each of the two eigenspaces. They are products of paired degree-two classes. The geometric endomorphism algebra is imaginary quadratic, so the Rosati calculation above gives N\'eron--Severi rank $n^2$, exactly the dimension of the paired degree-two space. The divisor cycle map fills that space.

For a supersingular elliptic curve the geometric endomorphism algebra is quaternionic and Rosati is its canonical conjugation. The self-adjoint matrices in $M_n(\operatorname{End}^0(E))$ have dimension
\[
 n+4\binom n2=n(2n-1)=\binom{2n}{2}.
\]
Thus divisor classes span all of $H^2$. The even exterior algebra on $H^1$ is generated by $H^2$, so divisor products span every even cohomology group over $\mathbb Q_\ell$.

To upgrade these spanning assertions to \emph{rational} algebraic classes, use Poincar\'e duality. It is perfect on the balanced monomial spaces in the ordinary case, and on all even cohomology in the supersingular case. Choose divisor-product bases $b_i,d_j$ of complementary spaces over $\mathbb Q_\ell$. Their intersection matrix $M=(\deg(b_i d_j))$ has rational entries and is invertible. If an algebraic cycle $z$ has
$\operatorname{cl}(z)=\sum_i c_i\operatorname{cl}(b_i)$, then its pairings with the $d_j$ give an invertible rational linear system with rational right-hand side. Hence all $c_i$ are rational. This proves the precise generation assertion.

Finally, \cite[Remark~3.7]{MilnePolarizations} proves that
$(-1)^r\deg(L^{n-2r}xy)$ is positive definite on primitive Lefschetz classes, for $0\le r\le n/2$, over every algebraically closed field. The class spaces just identified are all algebraic cohomology classes in the stated examples, giving the claimed instance of Hodge standard positivity.
\end{proof}

\begin{remark}
The complex argument uses only that $\operatorname{End}^0(E)$ is imaginary quadratic, so it also proves the same assertion for powers of any CM \emph{elliptic} curve. The finite-field argument applies to powers of every elliptic curve over a finite field. An arbitrary higher-genus curve with CM Jacobian, or an arbitrary CM abelian variety, is not covered by that argument. For general abelian varieties, \cite[Th.~3.3]{MilnePolarizations} is conditional on the Hodge conjecture for all complex CM abelian varieties; its unconditional Remark~3.7 concerns Lefschetz classes. No unspecified Masullo statement is substituted for that hypothesis.
\end{remark}

\begin{corollary}[SG-05: all abelian absolute values]
For every abelian variety $A/\mathbb F_q$, every eigenvalue on $H^i(A_{\overline{\mathbb F}_q},\mathbb Q_\ell)$ has absolute value $q^{i/2}$. This is available independently of the full Hodge standard conjecture in all codimensions.
\end{corollary}

\begin{proof}
The classical inputs are Rosati positivity, $\pi^\dagger\pi=q$, and the resulting degree-one bound, precisely \cite[Th.~13.2, Lemmas~14.2--14.3 and Th.~14.1]{MilneAV}. The exterior-algebra isomorphism \cite[Th.~11.1(b)]{MilneAV} makes every degree-$i$ eigenvalue a product of $i$ degree-one eigenvalues. Their moduli multiply to $q^{i/2}$. This is a short deduction from the cited theorem, not a new proof of Rosati positivity. The main sign-only theorem supplies a second proof of these absolute values.
\end{proof}

Hence the positivity mechanism already works in degree one for every abelian variety, and for all primitive algebraic classes of the elliptic powers just treated. None of these facts identifies the Gysin-remainder sign with an index form, or proves a new general positivity theorem for exotic algebraic classes.

\subsection{The letter's absolute values and all SA06 C.1 rows}

\begin{proposition}[SG-06: fibre purity and the invisible cycle kernel]
Every smooth projective finite-field fibre of a smooth projective spread of a complex variety in the supplied SA06 C.1 rows has the absolute values $q^{i/2}$ in degree $i$, by SG-08--SG-11 (equivalently, the known \cite[Th.~1.6]{WeilI}). These absolute values do not recover the kernel of the cycle class map or upgrade statements formulated modulo algebraic equivalence.
\end{proposition}

\begin{proof}
A projective embedding and finitely many equations descend to a finitely generated integral subring of $\mathbb C$. After shrinking its spectrum to the smooth locus one obtains a smooth projective model $\mathcal X\to S$. At each closed point of $S$ the residue field is finite, and the main sign-only theorem applies to the corresponding smooth projective fibre. This is meaningful even when the original very general complex example requires a transcendental field of definition; there is no canonical single finite-field Frobenius on that original complex fibre.

If $z$ has zero cohomology class and $w$ is a complementary algebraic cycle, compatibility of cycle class with products and trace gives
\[
 \deg(z\cdot w)=\operatorname{Tr}\bigl(\operatorname{cl}(z)\cup\operatorname{cl}(w)\bigr)=0.
\]
Every cohomological correspondence operator also factors through its cycle class. Consequently its traces, characteristic polynomials and local $L$-factors are unchanged on replacing a cycle by one with the same cohomology class. Spectral information therefore does not distinguish a nonzero Griffiths class from zero, or establish that it vanishes modulo algebraic equivalence. This is a factorization statement, not a claim of a finite-field counterexample to conjectured rational vanishing of Griffiths groups.
\end{proof}

The supplied intake provides full-check descriptions for C.1a--e and only abbreviated statement descriptions for C.1f--q. The table records every row without claiming to have independently checked the absent original proofs. In each row, ``purity'' may now be obtained by the reconstructed sign-only argument, not merely by citing Deligne. The remaining distinctions concern what purity does not assert about cycles.

{\small
\begin{longtable}{@{}p{0.08\textwidth}p{0.34\textwidth}p{0.52\textwidth}@{}}
\toprule
Row & Supplied variety or assertion & Consequence for this question \\
\midrule
\endfirsthead
\toprule
Row & Supplied variety or assertion & Consequence for this question \\
\midrule
\endhead
C.1a & Very general quintic threefold $V$; infinite Griffiths summand in $A^2_{\mathrm{alg}}$. & Every smooth finite-field fibre is pure. Its Griffiths kernel is invisible to traces. Its varying $H^3$ is not identified with the Tate line $\mathbb Q(-3)$; no number-field $H^3$-GRH follows. \\
C.1b & Failure of the letter's $A,A',A''$ modulo algebraic equivalence. & Fibre purity does not remove a cycle-class-kernel obstruction in that formulation. \\
C.1c & $L^{n-2}:\operatorname{NS}_{\mathbb Q}\to A^{n-1}_{\mathrm{alg}}$; examples $V\times\mathbb P^{n-3}$. & K\"unneth gives $H^i(V_s\times\mathbb P^m)=\bigoplus_j H^{i-2j}(V_s)(-j)$, with the required weights. The stated Griffiths cokernel remains separate. \\
C.1d & Endpoint cases, dimensions at most two, and surface Hodge index. & Curve purity also has the classical Hodge-index/Rosati proof. Parity is still not its index form. \\
C.1e & $V\times\mathbb P^1$; vacuous primitive letter-$B$ in the stated degree, with a nonzero numerical kernel. & Product purity does not imply algebraic equivalence equals numerical equivalence. The row's vacuity qualification is retained; this is not a counterexample to primitive Hodge standard positivity. \\
C.1f & Universal ranges of letter-$A$ over $\mathbb C$. & Purity in every cohomological degree does not enlarge valid ranges of a statement modulo algebraic equivalence. \\
C.1g & Universal quintic over $\mathbb Q(t_1,\ldots,t_{113})$ with a homologically trivial cycle. & Use a finite-type $\mathbb Z$-model and its smooth finite-field closed fibres. A generic transcendental base has no single arithmetic Frobenius. \\
C.1h & Geometric generic hyperplane and its Gysin image modulo algebraic equivalence. & Smooth closed fibres are pure; the asserted Chow-level Gysin image needs its separate proof. \\
C.1i & Universal $A'/A''$ range and split-quadric examples. & SG-02 gives the exact Tate eigenvalues for the quadrics. Their purity cannot repair an incorrectly quantified restriction or Gysin assertion. \\
C.1j & Middle letter-$B$ and injectivity on the Griffiths kernel. & Positivity on cohomological realizations and injectivity on an algebraic-equivalence kernel are different requirements. \\
C.1k & Cycle-class injectivity, finite rank or generation, and equivalence relations. & Frobenius purity neither supplies injectivity nor controls the size of its kernel. \\
C.1l & Radical correspondences constructed from Griffiths classes. & A zero cohomology class acts as zero on cohomology; its spectral data vanish without proving it zero modulo algebraic equivalence. \\
C.1m & Semisimplicity modulo algebraic equivalence in universal dimension ranges. & Purity alone does not give general Frobenius semisimplicity or semisimplicity of this correspondence algebra. Cote~16:049 expressly invokes $I(X\times X)$. \\
C.1n & An endpoint image criterion, standard $D$, and a kernel condition. & Weight purity supplies neither that cycle-equivalence assertion nor the extra kernel vanishing. \\
C.1o & Projective-space stabilization and a middle-positivity radical. & Tate shifts preserve the required weights; stabilization does not make a zero cycle class visible to Frobenius. \\
C.1p & Abelian projective stabilizations and Beauville/Griffiths pieces. & SG-05 and projective Tate shifts give purity. They do not determine the Beauville pieces of a Griffiths kernel. \\
C.1q & Blow-up exceptional normal-bundle polynomial and a radical. & Smooth projective finite-field blow-ups are pure; any exceptional-cycle or radical assertion remains a Chow-level assertion. The intake's truncated description does not authorize a stronger reconstruction. \\
\bottomrule
\end{longtable}
}

The answer to the letter's absolute-value question is therefore affirmative for every applicable finite-field fibre, using the sign-only proof. The stronger claim that this restores the letter's formulations modulo algebraic equivalence is not supported and is contradicted, in their stated complex scope, by the prior-run counterexamples being kept separate. The stipulated Dirichlet GRH transfers to the ambient Tate factor $\zeta(s-3)$ by SG-02. It does not establish a number-field GRH for an unshown non-Artin--Tate factor such as a general quintic's $H^3$.

\subsection{The ambient quadric versus the E.3u hypersurface}

\begin{proposition}[SG-07: restriction kills the ambient primitive line]
Let $j:Y\hookrightarrow Q^6$ be a smooth hypersurface of positive degree $d$, including the degree-$14$ section in the supplied SA06 E.3u description. The primitive class $\zeta\in H^6(Q^6,\mathbb Q_\ell)$ satisfies $j^*\zeta=0$.
\end{proposition}

\begin{proof}
The projection formula gives
\[
 j_*j^*\zeta=d\,h\zeta=0.
\]
The complement $U=Q^6\setminus Y$ is affine, since $Y$ is an effective ample divisor. Artin vanishing and Poincar\'e duality give $H_c^a(U,\mathbb Q_\ell)=0$ for $a<6$. The localization sequence therefore identifies $H^4(Q^6)\to H^4(Y)$. Dualizing gives an isomorphism
\[
 j_*:H^6(Y,\mathbb Q_\ell)
 \longrightarrow H^8(Q^6,\mathbb Q_\ell)(1).
\]
It follows that $j^*\zeta=0$. The twist is essential: both sides have Frobenius eigenvalue $q^3$ on the indicated Tate summand. The vanishing input is SGA~4, XIV, Th.~3.1; this is the ordinary ample-section proof of weak Lefschetz and requires no genericity assumption on the given smooth $Y$.
\end{proof}

This explains why the difference of the two plane traces can become a homologically trivial cycle on $Y$. Its morphic or Lawson nontriviality is additional structure. The single ambient Tate summand cannot by itself supply the variable $H^5(Y)$, nor see the Griffiths or Lawson kernel on $Y\times E$. The general sign-only theorem nevertheless proves the Frobenius absolute values of every smooth projective finite-field fibre of those varieties. The degree-$14$ nonvanishing assertion attributed to Walker and Nagel in SA06 E.3u is not independently verified by this argument, and is not needed for it.
